\topic{22 continued}{From a reading to a controlled experiment}
Read with a prediction in mind. Before running a case, write down the quantity
you expect to change, the quantities held fixed, and the evidence that would
make you reject your explanation. Then compute, compare and explain the result.

\textbf{Four complementary sources.}
\begin{enumerate}
\item \href{https://jose.theoj.org/papers/10.21105/jose.00021}{L. A. Barba and G. F. Forsyth (2018), \emph{CFD Python: the 12 steps to Navier--Stokes equations}, JOSE 1(9), 21.}
Use its progression from a worked example to a small parameter change and a
challenge. The finite-difference lessons include cavity and channel flow;
they are not a performance comparison of our three FV algorithms.
\item \href{https://www.kth.se/social/files/58ad81f0f276542156c566ab/project.pdf}{KTH, \emph{Computational Fluid Dynamics SG2212/SG3114}, Spring 2017, Task 1, Q2--Q6.}
Its MATLAB projection exercise connects stability, boundary conditions,
pressure-solver cost and mass balance. Q6 motivates the fixed-operator
experiment on the next pages. Our analytic trial fields and Python code are
new; they are not the KTH cavity trajectories.
\item \href{https://doi.org/10.1186/s40712-019-0104-7}{A. A. Boroujerdi and M. Hamzeh (2019), \emph{Comparison of varieties of numerical methods applied to lid-driven cavity flow}, IJMME 14, 7.}
Separate coupling from spatial discretization. Their algorithms are SIMPLE,
SIMPLEC and SIMPLER; their grid arrangements and first-/second-order upwind
schemes differ from this lesson. A ranking in their configurations is not a
universal Reynolds-number rule for SIMPLE/PISO/PIMPLE.
\item \href{https://doi.org/10.1002/(SICI)1097-0363(19980228)26:4<459::AID-FLD645>3.0.CO;2-U}{I. E. Barton (1998), \emph{Comparison of SIMPLE- and PISO-type algorithms for transient flows}, IJNMF 26, 459--483.}
Compare cost at a specified physical accuracy. That paper studies square-cylinder
and backward-facing-step flows and several time discretizations. Our startup
cavity keeps backward Euler fixed. We borrow the experimental question, not
the paper's geometries, second-order schemes or performance conclusions.
\end{enumerate}
\textbf{Use the existing evidence.} Tables 6a--6c and Figure 21 already separate
the mesh and convection effects. Explain why changing both a coupling algorithm
and an advection scheme would prevent a clean attribution of improved accuracy.
The new Tables L and T2 add linear-solver and physical-time experiments.
\newpage
\topic{22 continued}{Table L: pressure-solver work and conservation}
Freeze the positive-semidefinite pressure operator $A=-D d G$ with $d=0.05$ on a closed
unit square. $D$ maps face velocities to cell divergence; $G$ maps cell pressure
to face gradients. For each analytic trial velocity $q^*$, solve
\[
 b=-Dq^*,\qquad Ap=b,\qquad q=q^*-dGp.
\]
The normal wall velocity is zero, so $\sum_P b_P=0$. Remove one pressure degree
of freedom during the solve and subtract the mean afterward. The residual
of the \emph{full} operator, including the omitted gauge row, obeys
\[
 r=b-Ap,\qquad Dq=-r.
\]
Thus a loose linear solve can leave a mass defect even when the pressure
profile looks unchanged. Relative residual tolerances are not absolute
cell-divergence limits.
\begin{center}\small\begin{tabular}{rlrrr}\toprule
$N$ & Linear solve & Batch (ms) & Setup (ms) & $\|Dq\|_\infty$\\\midrule
16 & LU each RHS & 2.22 & 2.02 & 7.97e-14\\
16 & LU reused & 0.82 & 0.67 & 7.97e-14\\
16 & GMRES $10^{-4}$ & 113.09 & 0.23 & 6.28e-03\\
16 & GMRES $10^{-8}$ & 278.85 & 0.22 & 1.05e-06\\
32 & LU each RHS & 8.68 & 8.13 & 3.38e-12\\
32 & LU reused & 3.10 & 2.65 & 3.38e-12\\
32 & GMRES $10^{-4}$ & 818.02 & 0.24 & 3.83e-02\\
32 & GMRES $10^{-8}$ & 1807.16 & 0.22 & 4.21e-06\\
\bottomrule\end{tabular}\end{center}

Each batch has three different analytic right-hand sides at one fixed matrix.
Times are medians of three batches after an untimed warm-up on one Unity CPU.
Batch cost includes setup and solves; matrix assembly and diagnostics are outside
the timer. GMRES uses Jacobi preconditioning, restart 50 and a zero initial guess.
The setup and solve medians need not add to the median batch time.

\textbf{Exercise L.} Predict which cost LU reuse removes. Change the number of
right-hand sides from 3 to 6; report the batch cost and worst cell divergence
in one table. Tighten GMRES from $10^{-4}$ to $10^{-8}$ and compare \emph{both}
quantities. Record any nonzero solver return code instead of calling it a success.

\textbf{Reuse condition.} Reuse an LU factor only while every matrix entry is
unchanged. Changing the time step, grid, coefficients or boundary conditions
can invalidate it. In the cavity solver $d=1/\mathrm{diag}(A_{\rm momentum})$
generally changes between momentum assemblies; this fixed-Poisson experiment
does not license unconditional reuse inside SIMPLE or PIMPLE.
\newpage
\topic{22 continued}{Pressure agreement can hide a mass defect}
\lessonfigure{linear_pressure_profiles}{Fixed-operator projection on 32 cells per side, first analytic trial field, d=0.05. The upper panel compares zero-mean pressure on y/L=0.5. The lower panel shows remaining cell divergence along the bottom cell row, y/L=1/64, including the gauge-cell region. Its symmetric-log scale is linear below absolute divergence 1e-10. These are synthetic projection fields, not a cavity solution or a Ghia benchmark. LU and two Jacobi-preconditioned GMRES tolerances use the same matrix and right-hand side.}{linear-pressure-profiles}
\textbf{Explain the evidence.} Use the full residual identity to explain why
nearly overlapping pressure curves can coexist with different mass defects.
Distinguish the residual on the reduced linear system from conservation in
every pressure cell. Compare the largest defect with the declared acceptance
target, not only with the number of iterations.

\textbf{One-variable extension.} Keep the trial fields and grid fixed, change
only the GMRES tolerance, and estimate the tolerance needed for your chosen
absolute continuity target. A result on these small matrices is not evidence
that one linear solver is fastest on all CFD grids.
\newpage
\topic{22 continued}{Table T2: matched-time startup comparison}
All trajectories start from rest at $Re=100$ in the unit square with unit lid
speed, on the same 16-cell MAC mesh. Stop at $tU/L=0.2$. Hold central spatial
discretization and backward Euler fixed. PISO uses two inner corrections;
PIMPLE uses three outer passes and two inner corrections. Each trajectory is
repeated three times on one CPU; the table gives median solver-call wall time.
\begin{center}\small\begin{tabular}{lrrrr}\toprule
Method & $\Delta t$ & Velocity discrepancy (\%) & Seconds & Pressure solves\\\midrule
PISO & 0.05 & 6.259 & 0.021 & 8\\
PISO & 0.025 & 3.049 & 0.042 & 16\\
PISO & 0.0125 & 1.343 & 0.087 & 32\\
PIMPLE & 0.05 & 6.163 & 0.055 & 24\\
PIMPLE & 0.025 & 3.009 & 0.112 & 48\\
PIMPLE & 0.0125 & 1.320 & 0.218 & 96\\
\bottomrule\end{tabular}\end{center}
The measured comparison-field change after halving its time step is 0.4453\%; doubling its outer/inner correction counts changes it by 1.891e-14\%. These are sensitivity checks, not rigorous error bounds.

The comparison is PIMPLE at $\Delta t=0.003125$, six outer passes and four inner
corrections. It is checked against $\Delta t=0.00625$ and against twelve outer
passes/eight inner corrections at the fine step. For the native face-velocity
vector $q$, the table reports
\[
 E_t=100\,\frac{\|q_{\Delta t}(T)-q_{\rm comparison}(T)\|_2}
                         {\|q_{\rm comparison}(T)\|_2}.
\]
This is a discrepancy from a numerical comparison on the same mesh. It includes
time and splitting/nonlinear-coupling effects. It is not error against an exact
continuum solution. Steady Ghia profiles do not score a startup field.

The earlier Exercise B (Table T) uses cell-centered velocities and a different
comparison time step. Align the norm and comparison field before comparing
its percentages with Table T2. The earlier exercise remains available.

\textbf{Exercise T2.} First halve $\Delta t$ at fixed correction counts. Then
hold $\Delta t$ fixed and increase the correction counts. Compare profile
changes, the full nonlinear step defect and cost. Report whether the observed
reference sensitivity is small enough to support each ranking. More iterations
do not change backward Euler into a second-order time scheme.

\textbf{Report.} Submit Tables L and T2, two spatial profile comparisons,
the reference-sensitivity checks, and a short recommendation tied to a stated
accuracy target. Treat repeated timings as measurements on this host.
\newpage
\topic{22 continued}{Startup profiles at one physical time}
\lessonfigure{startup_time_profiles}{Unit-square cavity at Re=100, 16 MAC cells per side, tU/L=0.2, unit lid speed and initially stationary fluid. Upper panels show both centerline velocity components; lower panels show their signed differences from the fine PIMPLE comparison. Red is PISO and blue PIMPLE; line style denotes the time step. The black comparison uses dt=0.003125, six outer and four inner corrections. Spatial equations and backward Euler are common. The coarse mesh and impulsive lid startup limit accuracy claims.}{startup-time-profiles}
Locate the wall or return-flow regions where changing $\Delta t$ matters most.
Compare both velocity components before reducing the field to a single norm.
If reference sensitivity is comparable to a claimed improvement, refine the
comparison before selecting a winner. A low final continuity defect establishes
mass balance; it does not establish temporal accuracy or grid independence.
